\documentclass[preview]{standalone}

\usepackage{amsmath}
\usepackage{amssymb}
\usepackage{stellar}
\usepackage{definitions}
\usepackage{bettelini}

\begin{document}

\id{probability-finite-uniform-models}
\genpage

\section{Finite uniform models}

\begin{snippetdefinition}{finite-uniform-probability-definition}{Uniform probability on a finite sample space}
    Let \(\Omega\neq\emptyset\) be a finite \set.
    The \emph{uniform probability measure} on
    \((\Omega,\powerset(\Omega))\) is the \probabilitymeasure
    \[
        \mathbb P\colon\powerset(\Omega)\fromto[0,1],
        \qquad
        \mathbb P(A)=\frac{\cardinality{A}}{\cardinality{\Omega}}.
    \]
    Thus every \outcome \(\omega\in\Omega\) has the same probability
    \[
        \mathbb P(\{\omega\})=\frac{1}{\cardinality{\Omega}}.
    \]
    In a finite uniform model, computing probabilities is equivalent to
    counting favorable and possible outcomes correctly.
\end{snippetdefinition}

\begin{snippetexample}{introductory-counting-models-example}{Introductory counting models}
    The following finite uniform models illustrate the four basic counting
    patterns:
    \begin{enumerate}
        \item simplified Italian license plates with four letters and three
            digits are modeled by ordered choices with repetition, giving
            \(26^4\cdot 10^3\) possible plates;
        \item an ordered deck of \(40\) distinct cards is modeled by ordered
            choices without repetition, giving \(40!\) possible decks;
        \item a lottery draw of \(5\) numbers from \(90\), where order is not
            relevant and numbers are not repeated, gives
            \(\binom{90}{5}\) possible draws;
        \item distributing \(100\) indistinguishable coins among \(3\)
            distinguishable children, allowing a child to receive no coins,
            gives
            \[
                \binom{100+3-1}{100}=\binom{102}{2}
            \]
            possible distributions.
    \end{enumerate}
\end{snippetexample}

\section{Counting schemes}

\begin{snippetdefinition}{ordered-samples-with-replacement-definition}{Ordered samples with replacement}
    Let \(U\) be a finite \set with \(\cardinality{U}=n\), and let
    \(k\in\naturalnumbers\).
    The \emph{ordered samples with replacement} of length \(k\) from \(U\)
    are the \(k\)-tuples
    \[
        (x_1,\ldots,x_k)\in U^k.
    \]
    The same element of \(U\) may appear more than once, and the order of the
    tuple is part of the outcome.
\end{snippetdefinition}

\begin{snippetproposition}{ordered-samples-with-replacement-cardinality}{Number of ordered samples with replacement}
    Let \(U\) be a finite \set with \(\cardinality{U}=n\), and let
    \(k\in\naturalnumbers\). The number of \orderedwithreplacement of
    length \(k\) from \(U\) is
    \[
        n^k.
    \]
\end{snippetproposition}

\begin{snippetproof}{ordered-samples-with-replacement-cardinality-proof}{ordered-samples-with-replacement-cardinality}{Number of ordered samples with replacement}
    Each of the \(k\) positions can be chosen independently from \(n\) possible
    elements. By the fundamental principle of counting, the number of outcomes
    is
    \[
        \underbrace{n\cdot n\cdots n}_{k\text{ factors}}=n^k.
    \]
\end{snippetproof}

\begin{snippetdefinition}{ordered-samples-without-replacement-definition}{Ordered samples without replacement}
    Let \(U\) be a finite \set with \(\cardinality{U}=n\), and let
    \(k\in\naturalnumbers\) with \(k\leq n\).
    The \emph{ordered samples without replacement} of length \(k\) from \(U\)
    are the tuples
    \[
        (x_1,\ldots,x_k)\in U^k
        \suchthat
        x_i\neq x_j\quad\text{whenever }i\neq j.
    \]
    The order of the tuple is part of the outcome, and no element can be used
    twice.
\end{snippetdefinition}

\begin{snippetproposition}{ordered-samples-without-replacement-cardinality}{Number of ordered samples without replacement}
    Let \(U\) be a finite \set with \(\cardinality{U}=n\), and let
    \(k\in\naturalnumbers\) with \(k\leq n\). The number of
    \orderedwithoutreplacement of length \(k\) from \(U\) is
    \[
        n(n-1)\cdots(n-k+1)
        =
        \frac{n!}{(n-k)!}.
    \]
\end{snippetproposition}

\begin{snippetproof}{ordered-samples-without-replacement-cardinality-proof}{ordered-samples-without-replacement-cardinality}{Number of ordered samples without replacement}
    There are \(n\) choices for the first position, \(n-1\) choices for the
    second position, and so on. The \(k\)-th position has \(n-k+1\) choices.
    Multiplying gives
    \[
        n(n-1)\cdots(n-k+1)
        =
        \frac{n!}{(n-k)!}.
    \]
\end{snippetproof}

\begin{snippetdefinition}{unordered-samples-without-replacement-definition}{Unordered samples without replacement}
    Let \(U\) be a finite \set with \(\cardinality{U}=n\), and let
    \(k\in\naturalnumbers\) with \(k\leq n\).
    The \emph{unordered samples without replacement} of size \(k\) from \(U\)
    are the \set[subsets] \(A\subseteq U\) with \(\cardinality{A}=k\).
    The order of extraction is ignored, and no element can be used twice.
\end{snippetdefinition}

\begin{snippetproposition}{unordered-samples-without-replacement-cardinality}{Number of unordered samples without replacement}
    Let \(U\) be a finite \set with \(\cardinality{U}=n\), and let
    \(k\in\naturalnumbers\) with \(k\leq n\). The number of
    \unorderedwithoutreplacement of size \(k\) from \(U\) is
    \[
        \binom{n}{k}.
    \]
\end{snippetproposition}

\begin{snippetproof}{unordered-samples-without-replacement-cardinality-proof}{unordered-samples-without-replacement-cardinality}{Number of unordered samples without replacement}
    Every unordered sample of size \(k\) has exactly \(k!\) orderings, and
    these orderings are all distinct because there are no repetitions.
    Therefore the number of unordered samples is
    \[
        \frac{1}{k!}\frac{n!}{(n-k)!}
        =
        \binom{n}{k}.
    \]
\end{snippetproof}

\begin{snippetdefinition}{unordered-samples-with-replacement-definition}{Unordered samples with replacement}
    Let \(U=\{u_1,\ldots,u_n\}\) be a finite \set, and let
    \(k\in\naturalnumbers\).
    An \emph{unordered sample with replacement} of size \(k\) from \(U\) is
    specified by a vector
    \[
        (m_1,\ldots,m_n)\in\naturalnumbers^n
        \suchthat
        \sum_{i=1}^n m_i=k,
    \]
    where \(m_i\) is the number of times \(u_i\) appears.
    The order is ignored, and repetitions are allowed.
\end{snippetdefinition}

\begin{snippetproposition}{unordered-samples-with-replacement-cardinality}{Number of unordered samples with replacement}
    Let \(U\) be a finite \set with \(\cardinality{U}=n\), and let
    \(k\in\naturalnumbers\). The number of \unorderedwithreplacement
    of size \(k\) from \(U\) is
    \[
        \binom{n+k-1}{k}
        =
        \binom{n+k-1}{n-1}.
    \]
\end{snippetproposition}

\begin{snippetproof}{unordered-samples-with-replacement-cardinality-proof}{unordered-samples-with-replacement-cardinality}{Number of unordered samples with replacement}
    We must count the vectors
    \[
        (m_1,\ldots,m_n)\in\naturalnumbers^n
        \suchthat
        m_1+\cdots+m_n=k.
    \]
    Represent such a vector by \(k\) stars and \(n-1\) separators:
    \[
        \underbrace{\star\cdots\star}_{m_1}
        \mid
        \underbrace{\star\cdots\star}_{m_2}
        \mid
        \cdots
        \mid
        \underbrace{\star\cdots\star}_{m_n}.
    \]
    There are \(n+k-1\) positions in total, and choosing the \(k\) positions
    occupied by stars determines the sample. Hence the number is
    \[
        \binom{n+k-1}{k}.
    \]
\end{snippetproof}

\section{Particle systems}

\begin{snippetproposition}{maxwell-boltzmann-occupancy-model}{Distinguishable particles in finitely many states}
    Let \(n,r\in\naturalnumbers\). Consider \(n\) distinguishable particles,
    each of which can occupy one of \(r\) states. Assume every microstate in
    \[
        \Omega_{n,r}=\{1,\ldots,r\}^n
    \]
    has the same probability. Let
    \(k_1,\ldots,k_r\in\naturalnumbers\) satisfy
    \[
        k_1+\cdots+k_r=n.
    \]
    The probability that exactly \(k_i\) particles occupy state \(i\), for
    every \(i\), is
    \[
        \frac{n!}{r^n k_1!\cdots k_r!}.
    \]
\end{snippetproposition}

\begin{snippetproof}{maxwell-boltzmann-occupancy-model-proof}{maxwell-boltzmann-occupancy-model}{Distinguishable particles in finitely many states}
    The \samplespace \(\Omega_{n,r}\) has cardinality \(r^n\).
    For fixed occupation numbers \(k_1,\ldots,k_r\), the favorable microstates
    are obtained by choosing which \(k_i\) particles occupy state \(i\), for
    each \(i\). Their number is the multinomial coefficient
    \[
        \binom{n}{k_1,\ldots,k_r}
        =
        \frac{n!}{k_1!\cdots k_r!}.
    \]
    Dividing by \(\cardinality{\Omega_{n,r}}=r^n\) gives the formula.
\end{snippetproof}

\begin{snippetproposition}{bose-einstein-occupancy-model}{Indistinguishable particles in finitely many states}
    Let \(n,r\in\naturalnumbers\). Consider \(n\) indistinguishable particles
    that can occupy \(r\) states, with no restriction on how many particles can
    occupy the same state. The possible occupancy vectors are
    \[
        \widehat\Omega_{n,r}
        =
        \left\{
            (k_1,\ldots,k_r)\in\naturalnumbers^r
            \suchthat
            k_1+\cdots+k_r=n
        \right\}.
    \]
    If the \finiteuniformprobability is placed on \(\widehat\Omega_{n,r}\),
    then
    \[
        \cardinality{\widehat\Omega_{n,r}}
        =
        \binom{n+r-1}{n},
    \]
    and every occupancy vector has probability
    \[
        \binom{n+r-1}{n}^{-1}.
    \]
\end{snippetproposition}

\begin{snippetproof}{bose-einstein-occupancy-model-proof}{bose-einstein-occupancy-model}{Indistinguishable particles in finitely many states}
    The occupancy vectors are exactly the \unorderedwithreplacement
    of size \(n\) from the \set of \(r\) states.
    By the stars-and-bars formula,
    \[
        \cardinality{\widehat\Omega_{n,r}}
        =
        \binom{r+n-1}{n}.
    \]
    The stated probability follows from the \finiteuniformprobability on this
    finite \samplespace.
\end{snippetproof}

\section{Urn schemes}

\begin{snippetproposition}{red-blue-urn-with-replacement}{Red balls with replacement}
    Let \(N,R,n\in\naturalnumbers\) with \(R\leq N\). An urn contains \(N\)
    distinguishable balls, \(R\) of which are red and \(N-R\) of which are blue.
    Draw \(n\) balls with replacement, and suppose each ordered sequence of
    draws is equally likely. If \(A_k\) is the \event of drawing exactly
    \(k\) red balls, then
    \[
        \mathbb P(A_k)
        =
        \binom{n}{k}
        \left(\frac{R}{N}\right)^k
        \left(1-\frac{R}{N}\right)^{n-k}.
    \]
\end{snippetproposition}

\begin{snippetproof}{red-blue-urn-with-replacement-proof}{red-blue-urn-with-replacement}{Red balls with replacement}
    The \samplespace is the \set of \orderedwithreplacement of
    length \(n\), so it has cardinality \(N^n\).
    To draw exactly \(k\) red balls, choose the \(k\) positions of the red
    draws, choose one of the \(R\) red balls for each red position, and choose
    one of the \(N-R\) blue balls for each remaining position. Thus
    \[
        \cardinality{A_k}
        =
        \binom{n}{k}R^k(N-R)^{n-k}.
    \]
    Dividing by \(N^n\) gives the formula.
\end{snippetproof}

\begin{snippetproposition}{red-blue-urn-without-replacement}{Red balls without replacement}
    Let \(N,R,n\in\naturalnumbers\) with \(R\leq N\) and \(n\leq N\).
    An urn contains \(N\) distinguishable balls, \(R\) of which are red and
    \(N-R\) of which are blue. Draw \(n\) balls without replacement.
    If \(A_k\) is the \event of drawing exactly \(k\) red balls, then
    \[
        \mathbb P(A_k)
        =
        \frac{\binom{R}{k}\binom{N-R}{n-k}}{\binom{N}{n}},
    \]
    where the formula is understood for values of \(k\) such that
    \(0\leq k\leq R\) and \(0\leq n-k\leq N-R\).
\end{snippetproposition}

\begin{snippetproof}{red-blue-urn-without-replacement-proof}{red-blue-urn-without-replacement}{Red balls without replacement}
    Since only the selected balls matter, we may use the
    \finiteuniformprobability on the \unorderedwithoutreplacement of size
    \(n\).
    There are \(\binom{N}{n}\) possible samples.
    To obtain exactly \(k\) red balls, choose \(k\) balls among the \(R\) red
    balls and \(n-k\) balls among the \(N-R\) blue balls. Thus
    \[
        \cardinality{A_k}
        =
        \binom{R}{k}\binom{N-R}{n-k}.
    \]
    Dividing favorable outcomes by possible outcomes gives the formula.
\end{snippetproof}

\begin{snippetproposition}{hypergeometric-binomial-limit}{Sampling without replacement approaches sampling with replacement}
    Let \(n,k\in\naturalnumbers\) be fixed. Let \(N_m,R_m\in\naturalnumbers\)
    be such that \(R_m\leq N_m\), \(N_m\to+\infty\), and
    \[
        \frac{R_m}{N_m}\to p\in(0,1).
    \]
    Then
    \[
        \frac{\binom{R_m}{k}\binom{N_m-R_m}{n-k}}{\binom{N_m}{n}}
        \to
        \binom{n}{k}p^k(1-p)^{n-k}.
    \]
\end{snippetproposition}

\begin{snippetproof}{hypergeometric-binomial-limit-proof}{hypergeometric-binomial-limit}{Sampling without replacement approaches sampling with replacement}
    Write the ratio as
    \[
        \binom{n}{k}
        \frac{R_m(R_m-1)\cdots(R_m-k+1)}{N_m(N_m-1)\cdots(N_m-k+1)}
        \frac{(N_m-R_m)\cdots(N_m-R_m-(n-k)+1)}
        {(N_m-k)\cdots(N_m-n+1)}.
    \]
    For each fixed factor,
    \[
        \frac{R_m-a}{N_m-b}\to p
        \qquad\text{and}\qquad
        \frac{N_m-R_m-a}{N_m-b}\to 1-p.
    \]
    Since \(n\) and \(k\) are fixed, passing to the product gives the limit.
\end{snippetproof}

\begin{snippetproposition}{red-draw-position-probability}{Probability of a red ball at a fixed position}
    Let \(N,R,n\in\naturalnumbers\) with \(R\leq N\) and \(n\leq N\).
    An urn contains \(N\) distinguishable balls, \(R\) of which are red.
    Whether we draw with replacement or without replacement, if \(R_i\) is the
    \event that the \(i\)-th draw is red, then
    \[
        \mathbb P(R_i)=\frac{R}{N}
    \]
    for every \(i\in\{1,\ldots,n\}\).
\end{snippetproposition}

\begin{snippetproof}{red-draw-position-probability-proof}{red-draw-position-probability}{Probability of a red ball at a fixed position}
    With replacement, the statement follows immediately because each draw has
    \(R\) favorable balls out of \(N\).

    Without replacement, use the \finiteuniformprobability on
    \orderedwithoutreplacement. For each position \(i\), the map that swaps
    the first and \(i\)-th positions is a bijection of the \samplespace and
    sends the \event that the first draw is red to the \event that the
    \(i\)-th draw is red. Therefore all positions have the same probability
    of being red. Since
    the first draw is red with probability \(\frac{R}{N}\), so is every
    \(i\)-th draw.
\end{snippetproof}

\section{Repeated trials}

\begin{snippetdefinition}{finite-bernoulli-product-model}{Finite Bernoulli product model}
    Let \(N\in\naturalnumbers\) and \(p\in[0,1]\). The finite Bernoulli
    product model with \(N\) trials and success probability \(p\) is the
    \probabilityspace
    \[
        \Omega=\{0,1\}^N,
    \]
    where \(1\) represents success and \(0\) represents failure, and
    \[
        \mathbb P(\{\omega\})
        =
        p^{\sum_{i=1}^N\omega_i}
        (1-p)^{N-\sum_{i=1}^N\omega_i}
    \]
    for every \(\omega=(\omega_1,\ldots,\omega_N)\in\Omega\).
\end{snippetdefinition}

\begin{snippetproposition}{bernoulli-product-model-from-independence}{Finite Bernoulli product model from independent trials}
    Let \(N\in\naturalnumbers\) and \(p\in[0,1]\).
    Suppose \((\Omega,\powerset(\Omega),\mathbb P)\) is a
    \probabilityspace with \(\Omega=\{0,1\}^N\).
    For \(i\in\{1,\ldots,N\}\), let
    \[
        T_i=\{\omega\in\Omega\suchthat\omega_i=1\}.
    \]
    If \(\mathbb P(T_i)=p\) for every \(i\) and the \event[events]
    \(T_1,\ldots,T_N\) are \independentevents, then \(\mathbb P\) is the
    \bernoulliproductmodel.
\end{snippetproposition}

\begin{snippetproof}{bernoulli-product-model-from-independence-proof}{bernoulli-product-model-from-independence}{Finite Bernoulli product model from independent trials}
    Let \(C_i=\Omega\difference T_i\). By \eventindependence and the
    characterization of independence by complements, for every choice
    \(L_i\in\{T_i,C_i\}\),
    \[
        \mathbb P\left(\bigcap_{i=1}^N L_i\right)
        =
        \prod_{i=1}^N\mathbb P(L_i).
    \]
    For \(\omega\in\{0,1\}^N\), choose \(L_i=T_i\) if \(\omega_i=1\), and
    \(L_i=C_i\) if \(\omega_i=0\). Then
    \[
        \{\omega\}=\bigcap_{i=1}^N L_i,
    \]
    so
    \[
        \mathbb P(\{\omega\})
        =
        \prod_{i=1}^N p^{\omega_i}(1-p)^{1-\omega_i}
        =
        p^{\sum_{i=1}^N\omega_i}
        (1-p)^{N-\sum_{i=1}^N\omega_i}.
    \]
    A finite \probabilitymeasure is determined by its values on singletons,
    so this determines \(\mathbb P\).
\end{snippetproof}

\begin{snippetproposition}{binomial-count-from-bernoulli-trials}{Number of successes in finite Bernoulli trials}
    In the \bernoulliproductmodel with \(N\) trials and success
    probability \(p\), let \(S_k\) be the \event of obtaining exactly
    \(k\) successes. Then
    \[
        \mathbb P(S_k)
        =
        \binom{N}{k}p^k(1-p)^{N-k}.
    \]
\end{snippetproposition}

\begin{snippetproof}{binomial-count-from-bernoulli-trials-proof}{binomial-count-from-bernoulli-trials}{Number of successes in finite Bernoulli trials}
    Every \(\omega\in S_k\) has exactly \(k\) entries equal to \(1\), and
    hence
    \[
        \mathbb P(\{\omega\})=p^k(1-p)^{N-k}.
    \]
    There are \(\binom{N}{k}\) such \outcome[outcomes], one for each choice
    of the \(k\) success positions. Summing over \(S_k\) gives the formula.
\end{snippetproof}

\begin{snippetproposition}{first-success-finite-bernoulli-trials}{First success in finite Bernoulli trials}
    In the \bernoulliproductmodel with \(N\) trials and success
    probability \(p\), let \(B_n\) be the \event that the first success occurs
    on trial \(n\), where \(1\leq n\leq N\). Then
    \[
        \mathbb P(B_n)=p(1-p)^{n-1}.
    \]
\end{snippetproposition}

\begin{snippetproof}{first-success-finite-bernoulli-trials-proof}{first-success-finite-bernoulli-trials}{First success in finite Bernoulli trials}
    The \event \(B_n\) requires failures in the first \(n-1\) trials and a
    success at trial \(n\). The remaining \(N-n\) trials are unrestricted.
    Summing over all possibilities for the remaining trials gives
    \[
        \mathbb P(B_n)
        =
        (1-p)^{n-1}p
        \sum_{h=0}^{N-n}\binom{N-n}{h}p^h(1-p)^{N-n-h}.
    \]
    By the binomial theorem, the sum is \(1\), so
    \[
        \mathbb P(B_n)=p(1-p)^{n-1}.
    \]
\end{snippetproof}

\begin{snippetproposition}{at-least-one-success-finite-bernoulli-trials}{At least one success in finite Bernoulli trials}
    In the \bernoulliproductmodel with \(N\) trials and success
    probability \(p\), let \(D_n\) be the \event of obtaining at least one
    success in the first \(n\) trials, where \(1\leq n\leq N\). Then
    \[
        \mathbb P(D_n)=1-(1-p)^n.
    \]
\end{snippetproposition}

\begin{snippetproof}{at-least-one-success-finite-bernoulli-trials-proof}{at-least-one-success-finite-bernoulli-trials}{At least one success in finite Bernoulli trials}
    The complement of \(D_n\) is the \event that the first \(n\) trials are
    all failures. Its probability is \((1-p)^n\). Therefore
    \[
        \mathbb P(D_n)=1-\mathbb P(\Omega\difference D_n)=1-(1-p)^n.
    \]
\end{snippetproof}

\end{document}
